% Related Work - Communication-Efficient Federated Learning
% This file is input by the main document. Do not add \begin or \end.

Communication efficiency is a central challenge in federated learning,
where the cost of transmitting model updates between clients and the server
can dominate training time~\cite{kairouz2021fl}. Several compression
strategies have been proposed to mitigate this bottleneck.

Alistarh et al.~\cite{alistarh2017qsgd} introduced QSGD, which combines
stochastic gradient quantization with unbiased encoding to achieve convergence
guarantees under bounded variance. Bernstein et al.~\cite{bernstein2018signsgd}
proposed SignSGD, where only the sign of each gradient is communicated,
enabling 1-bit compression with majority-vote aggregation. Haddadpour et
al.~\cite{haddadpour2021fedcom} developed FedCOM, providing a unified analysis
of compressed FL with periodic communication and local gradient tracking.
Top-k sparsification~\cite{aji2017topk} further reduces communication by
transmitting only the largest gradient components. Hierarchical FL
approaches~\cite{liu2020hierarchical,abad2020hierarchical} introduce
multi-level aggregation to reduce coordinator load in large-scale deployments.

However, most prior work evaluates these techniques exclusively in simulation
environments with high-end hardware, leaving open questions about their
applicability to resource-constrained edge deployments. In this work, we
implement stochastic rounding quantization within a simulated FedAvg pipeline
and validate it on synthetic non-IID partitions, bridging the gap between
theoretical compression and practical edge deployment.
